\chapter{Introduction}%
\label{chp:introduction}

\section{Problem Definition}
The rapid advancement of text-to-image generative models has fundamentally transformed visual content creation, with diffusion-based architectures achieving unprecedented image quality and text-prompt controllability. Among these, Stable Diffusion \cite{Rombach_2022_CVPR} has emerged as a particularly influential model due to its open-source availability, widespread deployment across commercial and research applications, and the extensive ecosystem of tools and fine-tuned variants built upon it. 

These systems rely critically on massive web-scraped datasets to learn the mapping between textual descriptions and visual content. Re-LAION-5B, containing over 5 billion image-text pairs harvested from the internet, represents one of the largest publicly available training corpora for such models \cite{Schuhmann_2022_Laion5b}. However, the scale and automated nature of web scraping introduces systematic biases that reflect and potentially amplify existing societal inequities present in online content \cite{Birhane_2021_MultiModalDatasetsMisogynyPornography}.

While demographic bias in AI systems is well-documented across domains such as natural language processing \cite{Bolukbasi_2016_ManComputerProgrammerWoman}, machine learning systems \cite{Mehrabi_2022_SurveyBiasFairnessMachine}, and computer vision \cite{Buolamwini_2018_BiasGenderClassification}, there exists a gap between detecting bias in generative model outputs and understanding the specific low-level image features and dataset characteristics that cause and perpetuate these biases.
 
Current research predominantly adopts an outcome-focused perspective, measuring demographic distributions in generated images and comparing them against desired benchmarks \cite{Luccioni_2023_AnalyzingSocietalRepresentationsDiffusionModels, Bianchi_2023_TTIGenAmplifiesDemographicStereotypes}. This approach treats models as black boxes, diagnosing the symptoms of bias without investigating the underlying representational features that enable biased behavior.

A fundamental yet underexplored dimension of this problem involves spurious features—low-level, non-semantic image characteristics including color statistics, texture patterns, edge distributions, frequency domain properties, and compression artifacts \cite{Meister_2023_GenderArtifactsinVisualDatasets, Zeng_2024_UnderstandingBiasinLargeScaleVisualDatasets}. Unlike semantic features that capture meaningful visual content, spurious features represent superficial statistical regularities that can enable biased shortcuts across downstream tasks. Critically, when these spurious features correlate with sensitive demographic attributes, models can make predictions based on superficial cues rather than semantic understanding, creating an insidious form of bias that operates beneath standard evaluation metrics. Furthermore, it remains unclear how generative training processes transform these spurious patterns, making it difficult to predict how training data biases will manifest in deployed systems.

Existing bias mitigation approaches typically require expensive interventions: curating balanced training datasets \cite{Yang_2020_FairerDatasets}, implementing fairness-constrained optimization during training \cite{Zhang_2018_MitigatingUnwantedBiasesAdversarial}, or fine-tuning models on carefully constructed demographic distributions \cite{Friedrich_2023_FairDiffusionInstructingTextToImage}. These methods demand substantial computational resources and extensive human annotation, creating barriers for researchers and practitioners working with pre-trained models. Furthermore, outcome-focused mitigation strategies may successfully balance demographic distributions without addressing the underlying spurious features that enable biased shortcuts, potentially leading to superficial corrections that fail to generalize across contexts or evaluation metrics.

This thesis addresses these gaps through a mechanistic investigation of spurious features in the Re-LAION-5B dataset and Stable Diffusion v1.5 model. Rather than solely measuring demographic outcomes, this research applies established methodologies from prior work to expose and characterize the low-level image features that distinguish dataset sources and correlate with sensitive attributes. By understanding which specific spurious patterns enable biased behavior in this influential dataset-model pair, this work provides empirical evidence regarding bias mechanisms and evaluates targeted inference-time interventions that address root causes rather than symptoms.

\section{Motivation}

The main motivating factor is the critical role that spurious features within generative models play in the generative AI ecosystem and the urgent need for accessible fairness interventions.

\textbf{Influence and Impact:} Re-LAION-5B serves as training data for multiple widely-deployed text-to-image models, while Stable Diffusion v1.5 has achieved widespread adoption with millions of users generating billions of images through platforms like DreamStudio, Hugging Face, and numerous third-party integrations. Despite this pervasive influence, comprehensive analysis of the spurious features within Re-LAION-5B and their propagation through Stable Diffusion v1.5 remains limited. Understanding bias mechanisms in these specific systems directly impacts the fairness of deployed generative AI at scale.

\textbf{Computational Accessibility:} Training or fine-tuning large-scale diffusion models requires substantial computational resources placing such interventions beyond reach for most researchers and practitioners. By focusing on inference-time methods applicable to frozen models, this research demonstrates fairness interventions accessible to smaller research groups, independent developers, and organizations with limited computational budgets who cannot implement training-time solutions.

\textbf{Reproducibility and Transparency:} Much fairness research examines proprietary systems where training data and model weights remain inaccessible, limiting independent verification. The open-source nature of Stable Diffusion v1.5 and the research availability of Re-LAION-5B enable findings to be independently verified, replicated, and extended by the broader community, strengthening the credibility and impact of this work.\\

\textbf{Beyond Outcome Measurement:} While existing research predominantly measures demographic distributions in generated images, understanding \emph{why} biases emerge—specifically, which low-level features enable biased shortcuts—enables interventions that address root causes rather than symptoms. This mechanistic understanding is essential for developing robust fairness solutions that generalize across contexts rather than superficial corrections that fail when evaluation metrics or use cases change.


\section{Aims and Objectives}

The overarching aim of this thesis is to apply established spurious feature detection methodologies to the Re-LAION-5B dataset and Stable Diffusion v1.5 model, systematically characterise the spurious features present in each, quantify their correlation with sensitive demographic attributes, and evaluate the effectiveness of existing inference-time debiasing techniques when applied to this specific dataset-model pair.

To accomplish this aim, the research pursues the following specific objectives:

\textbf{Objective 1: Dataset Acquisition, Preparation, and Annotation}

Retrieve images from Re-LAION-5B and generate corresponding images using Stable Diffusion v1.5 across a pre-defined set of occupation categories to enable systematic bias analysis. This includes addressing technical challenges such as deriving image embeddings in the absence of pre-computed CLIP embeddings from the Re-LAION-5B dataset. Annotate all retrieved and generated images for gender and skin-tone attributes using established computer vision models to enable demographic analysis and spurious feature correlation assessment.

\textbf{Objective 2: Identify and Implement Spurious Feature Detection Methods}

Identify and implement established classifier-based methodologies for detecting spurious features that distinguish between dataset sources. Apply systematic image transformations (color extraction, edge maps, frequency filtering) to isolate non-semantic features and train discriminative classifiers to characterise dataset-specific spurious patterns in Re-LAION-5B versus Stable Diffusion v1.5 outputs.

\textbf{Objective 3: Assess Spurious Feature Correlations with Demographics}

Analyze whether identified spurious features exhibit systematic correlations with sensitive demographic attributes (gender and skin-tone) within both Re-LAION-5B and Stable Diffusion v1.5 generated images. Quantify these correlations to determine the extent to which spurious features enable demographic classification shortcuts, and compare demographic distributions across occupations against real-world statistics to identify representation biases.

\textbf{Objective 4: Identify and Implement Inference-Time Debiasing Methods}

Identify and implement established inference-time debiasing techniques applicable to Stable Diffusion v1.5 without model retraining. Adapt these methods based on identified spurious features and evaluate their effectiveness through multiple metrics including demographic representation balance, spurious feature mitigation, and maintenance of image quality. Assess trade-offs between fairness improvements and generation quality to inform practical deployment guidelines.

\textbf{Objective 5: Derive Conclusions on Bias Mechanisms and Mitigation Effectiveness}

Synthesize findings to characterize how spurious features manifest in Re-LAION-5B and propagate through Stable Diffusion v1.5, determining whether the generative process preserves, amplifies, or mitigates training data biases. Document the effectiveness of inference-time debiasing methods on a widely-deployed production model and provide evidence-based recommendations for practitioners seeking to deploy fairer text-to-image generation systems.

Through these objectives, this thesis provides empirical analysis of spurious features and demographic bias in Re-LAION-5B and Stable Diffusion v1.5, demonstrating the practical application of existing methodologies to real-world systems and contributing actionable insights for developing more equitable text-to-image generation technologies.

\section{Document Structure}

